\documentclass[11pt,a4paper]{article}
\usepackage[utf8]{inputenc}
\usepackage[T1]{fontenc}
\usepackage{lmodern}
\usepackage{amsmath,amssymb}
\usepackage{booktabs}
\usepackage{xcolor}
\usepackage{geometry}
\usepackage[colorlinks=true,linkcolor=blue!50!black,urlcolor=blue!50!black]{hyperref}
\geometry{a4paper,left=2.2cm,right=2.2cm,top=2.2cm,bottom=2.4cm}
\setlength{\parindent}{0pt}
\setlength{\parskip}{0.7ex}
\numberwithin{equation}{section}
\newcommand{\code}[1]{\texttt{\small #1}}
\newcommand{\dg}{^{\dagger}}

\title{Several GPUs for domain-wall fermions in RutticeQCD}
\author{Wei-Lun Chen}
\date{\today\quad(internal note)}

\begin{document}
\maketitle
\tableofcontents

%======================================================================
\section{Which data a GPU needs from its neighbor}

The hopping term at site $x$ is
\begin{equation}
(H\psi)(x) = -\kappa\sum_{\mu=0}^{3}\Big[
\underbrace{(1-s\gamma_\mu)U_\mu(x)\,\psi(x+\hat\mu)}_{\text{forward leg}}
+\underbrace{(1+s\gamma_\mu)U_\mu(x-\hat\mu)\dg\,\psi(x-\hat\mu)}_{\text{backward leg}}\Big],
\end{equation}
with $s=+1$ for $H$ and $s=-1$ for $H\dg$.

A GPU holds the sites $0\le x_\mu<L_\mu$. For a site with $x_\mu=L_\mu-1$ the forward leg needs
$\psi(x+\hat\mu)$, which is on the next GPU. For $x_\mu=0$ the backward leg needs
$\psi(x-\hat\mu)$ and the link $U_\mu(x-\hat\mu)$, both on the previous GPU.

\subsection{Only two spin components travel}

Take $\mu=0$. In the Dirac representation used by the kernels,
\begin{equation}
\gamma_0=\begin{pmatrix}0&0&0&-i\\0&0&-i&0\\0&i&0&0\\i&0&0&0\end{pmatrix}.
\end{equation}
Apply $1-s\gamma_0$ to $w=(w_0,w_1,w_2,w_3)$:
\begin{align}
\big[(1-s\gamma_0)w\big]_0 &= w_0 - s(-i)w_3 = w_0 + is\,w_3,\\
\big[(1-s\gamma_0)w\big]_1 &= w_1 - s(-i)w_2 = w_1 + is\,w_2,\\
\big[(1-s\gamma_0)w\big]_2 &= w_2 - s(i)w_1 = w_2 - is\,w_1,\\
\big[(1-s\gamma_0)w\big]_3 &= w_3 - s(i)w_0 = w_3 - is\,w_0.
\end{align}
Compare row 2 with row 1:
\begin{align}
-is\,\big[(1-s\gamma_0)w\big]_1 &= -is\,w_1 - is\cdot is\,w_2\\
&= -is\,w_1 - i^2 s^2\,w_2\\
&= -is\,w_1 + w_2 \qquad (i^2=-1,\ s^2=1)\\
&= \big[(1-s\gamma_0)w\big]_2 .
\end{align}
In the same way row 3 $= -is\times$ row 0. So rows 2 and 3 follow from rows 0 and 1, and we
send only
\begin{equation}
h_a = \big[(1-s\gamma_\mu)w\big]_a,\qquad a=0,1,
\end{equation}
two color vectors: $2\times3$ complex numbers $=12$ reals instead of 24. The receiver rebuilds
rows 2 and 3. For $\gamma_3=\mathrm{diag}(1,1,-1,-1)$ it is even simpler:
\begin{equation}
(1-\gamma_3)w = (0,\,0,\,2w_2,\,2w_3),\qquad (1+\gamma_3)w = (2w_0,\,2w_1,\,0,\,0).
\end{equation}

\subsection{The sender applies the backward link}

The backward leg needs $U_\mu(x-\hat\mu)$, a link of the sending GPU. So the sender computes
\begin{equation}
t_a = U_\mu(y)\dg\,\big[(1+s\gamma_\mu)w(y)\big]_a,\qquad y = x-\hat\mu,
\end{equation}
and the receiver uses $t_a$ as it is. For the forward leg the sender sends
\begin{equation}
h_a = \big[(1-s\gamma_\mu)w(y)\big]_a,\qquad y=x+\hat\mu,
\end{equation}
and the receiver multiplies by its own $U_\mu(x)$. No GPU needs a link of another.

\subsection{The fifth dimension is applied before sending}

For domain-wall fermions the hop acts on
\begin{align}
w_s &= b\,\psi_s - c\,(Y\psi)_s,\\
(Y\psi)_s &= -P_-\psi_{s+1} - P_+\psi_{s-1}\qquad(+m\text{ instead of }-1\text{ across the walls}).
\end{align}
$w_s$ needs $\psi_{s\pm1}$ of the same four-dimensional site, which the sender has. So the sender
computes $w$ at its face sites, then projects. All $L_s$ values of $s$ of a site go in one
message. Bridge++ does the same (its \code{5dir} before \code{hop1}).

%======================================================================
\section{Why the result is the same bit for bit}

\subsection{Floating-point addition is not associative}

In double precision:
\begin{align}
(10^{16} + 1) - 10^{16} &= 10^{16} - 10^{16} = 0 \qquad(10^{16}+1\text{ rounds to }10^{16}),\\
10^{16} + (1 - 10^{16}) &= 10^{16} - (10^{16}-1) = 1 .
\end{align}
The same three numbers give 0 or 1. Usually only the last bits differ, but they do differ.

\subsection{The order of the eight legs}

One GPU sums the legs $\ell_0,\dots,\ell_7$ of a site in a fixed order:
\begin{equation}
S_{\text{one}} = \big(\big(\big(\ell_0+\ell_1\big)+\ell_2\big)+\cdots\big)+\ell_7 .
\end{equation}
Bridge++ on several GPUs sums the local legs first (bulk kernel) and adds the remote ones later
(boundary kernel). For a site whose legs 1 and 5 are remote:
\begin{equation}
S_{\text{Bridge++}} = \big(\ell_0+\ell_2+\ell_3+\ell_4+\ell_6+\ell_7\big) + \big(\ell_1+\ell_5\big)
\ \neq\ S_{\text{one}}\quad\text{in the last bits.}
\end{equation}
We compute such a site entirely in the surface kernel, in the order of $S_{\text{one}}$, reading
$\ell_1,\ell_5$ from the received buffer:
\begin{equation}
S_{\text{ours}} = \big(\big(\big(\ell_0+\ell_1^{\text{buf}}\big)+\ell_2\big)+\cdots\big)+\ell_7 .
\end{equation}
Each leg is the same arithmetic in both cases (projection on the sender, link product, rebuild on
the receiver), so $\ell_1^{\text{buf}}=\ell_1$ exactly and $S_{\text{ours}}=S_{\text{one}}$
exactly. The interior sites (no remote leg) are computed while the faces travel; no site is
computed twice.

\subsection{The same arithmetic: fused multiply-adds}

``The same arithmetic'' needs one more thing. A GPU has a fused multiply-add,
$\mathrm{fma}(a,b,c) = \mathrm{round}(ab+c)$, one rounding instead of two:
\begin{equation}
\mathrm{round}\big(\mathrm{round}(ab)+c\big)\ \neq\ \mathrm{round}(ab+c)\quad\text{in the last bit, in general.}
\end{equation}
nvcc turns $ab+c$ into an fma where it likes (\code{-fmad}, on by default), and it decides kernel by
kernel. A leg is computed in different kernels: the sender's packing kernel, the receiver's
surface kernel, the one-process kernel. When the Wilson operator was given its own instance of
the hop kernel (constants $L_s=1$, $b=1$, $c=0$), nvcc fused the link product differently there
than in the packing kernel, and some surface sites differed from one process by $10^{-16}$. So
the sums of a leg are written as explicit fmas in a fixed order (\code{c\_macc}, \code{add\_leg},
the third row of a link, the clover blocks), and a product that must stay unfused is
\code{\_\_dmul\_rn} (\code{mul\_rn}), which nvcc never fuses:
\begin{align}
c &\leftarrow \mathrm{fma}(a_r,b_r,c_r),\ \mathrm{fma}(-a_i,b_i,\cdot)\quad\text{(real part of } c+ab),\\
c &\leftarrow \mathrm{fma}(a_r,b_i,c_i),\ \mathrm{fma}(a_i,b_r,\cdot)\quad\text{(imaginary part)}.
\end{align}
Every kernel then rounds a leg the same way, whatever nvcc would have chosen.

%======================================================================
\section{How a face is numbered}

\subsection{The site number of the device fields}

A site is numbered lexically, $x$ fastest:
\begin{equation}
n(x,y,z,t) = x + L_x\big(y + L_y(z + L_z\,t)\big).
\label{eq:lex}
\end{equation}
For the face $t=L_t-1$:
\begin{equation}
n = \underbrace{L_xL_yL_z(L_t-1)}_{\text{offset}} + \big(x+L_x(y+L_yz)\big),
\end{equation}
one contiguous range. For the face $x=L_x-1$:
\begin{equation}
n = (L_x-1) + L_x\big(y + L_y(z + L_z\,t)\big),
\end{equation}
so consecutive $y$ give sites $L_x$ apart. Three of the four faces are not contiguous, and a
kernel must gather them.

\subsection{The face number}

Let $d_0<d_1<d_2$ be the directions other than $\mu$. We number a site of the face by
\begin{equation}
f = c_{d_0} + L_{d_0}\big(c_{d_1} + L_{d_1}\,c_{d_2}\big).
\label{eq:face}
\end{equation}
Two neighbors across the face differ only in $c_\mu$:
\begin{equation}
x=(\dots,c_\mu=L_\mu-1,\dots)\ \text{on GPU }A,\qquad x+\hat\mu=(\dots,c_\mu=0,\dots)\ \text{on GPU }B,
\end{equation}
with the same $c_{d_0},c_{d_1},c_{d_2}$. Hence $A$ and $B$ compute the same $f$ from
Eq.~\eqref{eq:face}. Nothing else has to be agreed on: no table is kept or sent.

\subsection{On the sites of one parity}

The parity of a site is
\begin{equation}
p = (c_\mu + c_{d_0} + c_{d_1} + c_{d_2}) \bmod 2 .
\end{equation}
For fixed $c_\mu,c_{d_1},c_{d_2}$, solve for $c_{d_0}$:
\begin{align}
c_{d_0} &\equiv p - (c_\mu+c_{d_1}+c_{d_2}) \pmod 2,\\
c_{d_0} &= 2k + r,\qquad r = (p + c_\mu+c_{d_1}+c_{d_2})\bmod 2,\qquad k=0,\dots,\tfrac{L_{d_0}}{2}-1 .
\end{align}
So exactly one site of each pair $(2k,2k+1)$ has the parity, and $k=\lfloor c_{d_0}/2\rfloor$
numbers it:
\begin{equation}
f = \lfloor c_{d_0}/2\rfloor + \tfrac{L_{d_0}}{2}\big(c_{d_1} + L_{d_1}\,c_{d_2}\big).
\end{equation}
This needs $L_{d_0}$ even. We require all local extents even. Then the local parity is also the
global one:
\begin{equation}
\sum_\mu (\text{origin}_\mu + c_\mu) \equiv \sum_\mu c_\mu \pmod 2,\qquad
\text{origin}_\mu = k_\mu L_\mu\ \text{even.}
\end{equation}

\subsection{Order in the buffer, and coalescing}

The buffer of a face holds, for vector $v$, fifth coordinate $s$, real $r$ ($0\le r<12$) and face
site $f$,
\begin{equation}
\text{index}(v,s,r,f) = \big((v\,L_s + s)\cdot12 + r\big)\cdot N_f + f .
\end{equation}
In the packing kernel thread $j$ of a warp has $f=f_0+j$. For fixed $(v,s,r)$:
\begin{equation}
\text{index}(f_0+j) = \text{index}(f_0) + j,\qquad j=0,\dots,31:
\end{equation}
32 consecutive reals, one memory transaction for the warp.

The reads from the field follow the site number. A field has component $c$ of site $n$ at
\begin{equation}
a(c,n) = \big(\lfloor n/32\rfloor\,N_c + c\big)\cdot32 + (n\bmod 32).
\end{equation}
\begin{itemize}
\item Faces in $y,z,t$: $d_0=x$, so $f_0+j$ has $x_0+j$ and by Eq.~\eqref{eq:lex} $n=n_0+j$;
within a block of 32, $a(c,n_0+j)=a(c,n_0)+j$: coalesced.
\item Face in $x$: $d_0=y$, so $f_0+j$ has $y_0+j$ and $n=n_0+jL_x$; the addresses are $L_x$
apart: not coalesced.
\end{itemize}
Only the face in $x$ is read with a stride. It has $V/L_x$ sites like the others, and packing is
small next to the interior, so we left it.

The fused hop-and-matrix kernel uses an XOR on the column of a tile in shared memory,
\code{column \^{} (8*(row \& 1))}, against bank conflicts. The face buffers are in global memory
and read along $f$, so no XOR is needed there.

%======================================================================
\section{How the faces travel}

\subsection{The sequence}

For one application of the hops:
\begin{align*}
&1.\ \text{pack every cut face (default stream)}\\
&2.\ \text{send and receive the faces}\\
&3.\ \text{interior kernel, at the same time as 2}\\
&4.\ \text{wait for 2}\\
&5.\ \text{surface kernel}
\end{align*}
\textbf{NCCL}: step 2 runs on a stream of its own; an event after step 1 starts it, an event after
it makes step 5 wait, so the host never blocks.

\textbf{Through the host}: step 2 is copy to the host, MPI exchange, copy back; the interior
kernel runs while MPI exchanges.

\textbf{Choice}: NCCL if every process can make its communicator:
\begin{equation}
\text{failed} = \sum_{\text{processes}} \big[\text{NCCL failed here}\big],\qquad
\text{use NCCL}\iff \text{failed}=0,
\end{equation}
so all processes choose the same.

\textbf{Two messages between one pair.} If a direction is cut in two, GPU $A$ sends GPU $B$ its
face forward and its face backward. NCCL matches the messages of a pair in the order they are
posted. Posting the sends in the order of the faces $f=0,\dots,7$ and the receives in the order
of the faces the data left by keeps the two orders the same on both sides.

\subsection{NVSHMEM: an evaluation (not implemented)}

\textbf{A model of one application.} Let an operator take $t_c$ to compute on a GPU and an exchange
$t_x = t_0 + b/B$: a fixed latency $t_0$ and $b$ bytes at bandwidth $B$. With the interior
overlapping the exchange,
\begin{equation}
t_{\text{apply}} \approx t_{\text{surface}} + \max\big(t_{\text{interior}},\ t_x\big).
\end{equation}
Two limits follow:
\begin{align}
t_{\text{interior}} \gg t_x &:\quad t_{\text{apply}}\approx t_c\quad(\text{the exchange is hidden; its kind does not matter}),\\
t_{\text{interior}} \ll t_x &:\quad t_{\text{apply}}\approx t_x\approx t_0\quad(\text{the latency is all}).
\end{align}

\textbf{Where each case is.} The fine operator on a large local lattice is the first case. A local
lattice $16^4$ with $L_s=16$ has
\begin{align}
\text{interior work} &\propto 16^4\cdot16 \approx 10^6\ \text{site-slices, about a millisecond},\\
\text{a face} &= 16^3\cdot16\cdot12\ \text{reals}\cdot8\ \text{B}\approx 6\ \text{MB, about }0.1\ \text{ms on NVLink},
\end{align}
so the exchange hides. The coarse lattice is the second case. With blocks of $4^4$ the coarse local
lattice is $4^4=256$ sites; a face is 64 sites of $\text{dim}\approx 40$ complex numbers,
\begin{equation}
64\cdot40\cdot16\ \text{B}\approx 40\ \text{kB},\qquad t_c \sim 10\ \mu\text{s},
\end{equation}
and the coarse solve applies the operator hundreds of times per fine iteration.

\textbf{Latencies (typical, not measured here).} An NCCL group of sends and receives costs a launch
and its protocol, $t_0\sim10$--$20\ \mu$s; an NVSHMEM put with a signal within a node a few $\mu$s.
On the coarse lattice then
\begin{equation}
\frac{t_{\text{apply}}^{\text{NCCL}}}{t_{\text{apply}}^{\text{NVSHMEM}}} \approx \frac{t_c + 15\ \mu\text{s}}{t_c + 4\ \mu\text{s}} \approx \frac{25}{14} \approx 1.8 ,
\end{equation}
a 1.5--3$\times$ faster coarse solve depending on $t_c$; on the fine operator a few per cent at most.

\textbf{Why NVSHMEM can be faster there.} (i) No group launch: the put is issued on the stream, or
by the packing kernel itself. (ii) Within a node the packing kernel can write straight into the
neighbor's buffer through a peer pointer (\code{nvshmem\_ptr}): the face is never copied a second
time, and our layout ($f$ fastest) already makes those stores coalesced. (iii) The wait is a signal
on the stream (or in the kernel), not a host round trip. (iv) NCCL's kernels occupy multiprocessors
while they move data; NVSHMEM's puts use the copy engines or the network.

\textbf{What it would take.} Receive buffers and signals on the symmetric heap (one collective
allocation, sized for the largest number of vectors, doubled so a GPU one step ahead never
overwrites what its neighbor reads); a feature \code{nvshmem}; bootstrap by the same broadcast of a
unique identifier as NCCL; first the stream-ordered form
\code{putmem\_signal\_on\_stream}/\code{signal\_wait\_until\_on\_stream}, then the direct stores within
a node. NCCL and NVSHMEM can live in one program (NCCL for the fine operator and the sums, NVSHMEM
for the coarse ghosts); the symmetric heap should then be kept small
(\code{NVSHMEM\_SYMMETRIC\_SIZE}). Recent NCCL versions add symmetric memory and device-initiated
operations of their own, which might give the same without NVSHMEM; that has to be checked against
their documentation first.

\textbf{When to do it.} Only on a machine with several GPUs, where it can be measured: our machine
has one, so neither path can be timed here.

%======================================================================
\section{The gauge field: an extended lattice}

A staple at $x$ in direction $\mu$, plane $\nu$, reads
\begin{equation}
U_\nu(x)\,U_\mu(x+\hat\nu)\,U_\nu(x+\hat\mu)\dg + U_\nu(x-\hat\nu)\dg\,U_\mu(x-\hat\nu)\,U_\nu(x-\hat\nu+\hat\mu).
\end{equation}
The last link sits at $x-\hat\nu+\hat\mu$, a diagonal neighbor. A rectangle reaches
$x+2\hat\mu+\hat\nu$.

We copy the links into an extended lattice with a halo of depth $d$ ($d=1$ staples, $d=2$
rectangles):
\begin{equation}
-d \le c_\mu < L_\mu + d\quad\text{in every cut direction.}
\end{equation}
The halo is filled one direction after the other. With two cut directions, $x$ then $y$:
\begin{align*}
&\text{after }x:\quad \text{filled}=\{-d\le c_x<L_x+d,\ 0\le c_y<L_y\},\\
&\text{send in }y\text{ the layers }0\le c_y<d\text{ and }L_y-d\le c_y<L_y,\ \text{for all }-d\le c_x<L_x+d,\\
&\text{after }y:\quad \text{filled}=\{-d\le c_x<L_x+d,\ -d\le c_y<L_y+d\}.
\end{align*}
The corners $(c_x<0,\ c_y<0)$ came with the $y$ layers, which already carried the $x$ halo; no
diagonal message is needed. This is the CPU's \code{gauge::paths::Extended} on the device.

The kernels compute only the local sites and read the extended field; every path from a local site
stays inside it. The arithmetic of a staple is unchanged, so the result is that of one GPU bit for
bit. The fermion force uses the face exchange of the hops with $b=1$, $c=0$, i.e.\ $w=\psi$.

%======================================================================
\section{The coarse lattice}

\subsection{Aggregates inside a process}

With a block size $B_\mu$ dividing the local extent $L_\mu$:
\begin{equation}
L_\mu = k_\mu B_\mu \quad\Rightarrow\quad \text{blocks } \big[jB_\mu,(j+1)B_\mu\big),\ j=0,\dots,k_\mu-1,
\end{equation}
none crosses $c_\mu=L_\mu$. Restriction and prolongation stay local.

\subsection{Probing across a cut}

The link from coarse site $X$ to its neighbor in direction $\mu$ is
\begin{equation}
B_\mu(X)_{ij} = \big[P\dg\,H_\mu\,P\,e_j\big]_i(X),
\end{equation}
with $H_\mu$ the forward leg alone and $e_j$ the unit vector of component $j$ on every coarse site.
$H_\mu$ at a fine site of $X$ on the face reads $Pe_j$ on the next GPU: a face exchange of the field
$Pe_j$, done once per batch of $j$ and used for all eight legs.

\subsection{Ghost sites}

The coarse kernel reads a neighbor through a table $N[d][X]$. On several processes it points past
the local sites:
\begin{equation}
N[d][X] = \begin{cases}
\text{local neighbor}, & \text{inside the process},\\
n_{\text{sites}} + \text{first}[d] + f(X), & \text{across a cut},
\end{cases}
\qquad \text{first}[d]=\sum_{d'<d}\text{count}[d'],
\end{equation}
with $f$ the face number of Eq.~\eqref{eq:face} on the coarse lattice. Before each application the
faces of the input are gathered, exchanged and placed at those indices of a separate ghost field.
The vectors keep only local sites, so an inner product is a local sum followed by a sum over
processes.

\subsection{The exact adjoint}

For $(Av)(X)=A(X)v(X)+\sum_d B_d(X)\,v(N[d][X])$:
\begin{align}
\langle u, Av\rangle &= \sum_X u(X)\dg\Big[A(X)v(X) + \sum_d B_d(X)v(N[d][X])\Big]\\
&= \sum_X \big[A(X)\dg u(X)\big]\dg v(X) + \sum_d\sum_X \big[B_d(X)\dg u(X)\big]\dg v\big(N[d][X]\big).
\end{align}
In the second term put $Y=N[d][X]$. Then $X$ is the site whose hop $d$ comes from $Y$, i.e.\
$X=N[\bar d][Y]$ with $\bar d$ the opposite direction:
\begin{equation}
\langle u, Av\rangle = \sum_Y \Big[A(Y)\dg u(Y) + \sum_d B_d\big(N[\bar d][Y]\big)\dg u\big(N[\bar d][Y]\big)\Big]\dg v(Y),
\end{equation}
so
\begin{equation}
(A\dg u)(Y) = A(Y)\dg u(Y) + \sum_d B_d\big(N[\bar d][Y]\big)\dg\,u\big(N[\bar d][Y]\big).
\end{equation}
When $N[\bar d][Y]$ is a ghost its link $B_d$ lives on the neighboring process. We exchange the
links of the face sites once (at setup) and form the conjugate transposes; a conjugate transpose
is a copy with signs, so the numbers are those of one GPU.

\subsection{The levels below}

Each level is a lattice cut on the same process grid; a super-block of $b$ sites per direction must
divide the local extent of the level above, so it stays inside a process.

Its operator is probed with colors. With $k_\mu$ the smallest divisor $\ge3$ of the global extent,
\begin{equation}
\text{color}(c) = \sum_\mu \big(c_\mu \bmod k_\mu\big)\prod_{\nu<\mu}k_\nu,
\end{equation}
of the \emph{global} coordinates $c$. A ghost stands for a site of the neighbor whose global
coordinates we know, so it gets that site's color, and probing finds the links across the cut.

The last level is solved with its dense inverse. Every process gathers the whole level: each puts
its sites at their global places in an array of zeros, and the arrays are summed over processes,
\begin{equation}
\text{whole}[g] = \sum_{p} \text{part}_p[g],\qquad
\text{part}_p[g] = \begin{cases}\text{value}, & g\text{ on }p,\\ 0, & \text{otherwise},\end{cases}
\end{equation}
and $\text{value}+0+\dots+0=\text{value}$ exactly. Every GPU then inverts the same matrix. (The CPU
does the same for its levels.)

%======================================================================
\section{Sums over the processes}

An inner product on $P$ processes is
\begin{equation}
\langle a,b\rangle = \sum_{p=1}^{P}\Big(\sum_{x\in p} a(x)^* b(x)\Big),
\end{equation}
another order of the same terms than one sum over all $x$. By Sec.~2.1 the last bits differ, so a
solver is not bit for bit. Example: CGNE to $|r|^2/|b|^2<10^{-22}$ stopped after 1247 iterations on
two and four processes and 1246 on one; the solutions differ by $10^{-11}$ relative. The operators,
which have no such sum, are bit for bit.

%======================================================================
\section{What was checked}

Tests run several processes as threads of one program, all on the one A30, against one process:
\begin{center}
\begin{tabular}{ll}
\toprule
what & result \\
\midrule
$D$, $D\dg$, $\Gamma_5D$, $\tilde A$ (both forms, $\alpha=1.3$), up to 16 ranks & bit for bit \\
Wilson and clover $D$, $D\dg$, $\gamma_5D$, hops, $T$, $T^{-1}$ (up to 16 ranks) & bit for bit \\
their HMC derivatives (hops, clover leaves, $\ln\det T_o$) & bit for bit \\
stout (2 steps), chain rule, rectangle force, fermion force & bit for bit \\
P, coarse operator, $D_c$, $\Gamma_5D_c\Gamma_5$, $D_c\dg$, Schur & bit for bit \\
levels: quark V-cycle, $|H|^{-1}$ of the Pauli-Villars factor & bit for bit \\
Wilson/clover multigrid: P, test vectors (up to $2\times2\times2\times2$) & bit for bit \\
Wilson/clover multigrid: FGMRES with it & same iterations, $10^{-15}$ \\
CGNE (domain-wall; clover Schur) & iterations within 0.2\%, $10^{-11}$ \\
\bottomrule
\end{tabular}
\end{center}
Every bit-for-bit test was also broken on purpose (wrong boundary factor, halo too shallow, ghosts
on the wrong face, wrong ghost colors, ghost links of the wrong direction, the surface of the
Wilson and clover operators skipped) and then failed.

With two MPI processes on $4^3\times8$, GPU and CPU alike: the propagator (CG, mixed CG, the
multigrid with and without a second level, MINRES or GCR for the Pauli-Villars factor) and HMC
(CG, mixed CG, the multigrid with a second level) take the same iterations as one process,
$\Delta H$ to $10^{-10}$.

We have not checked NCCL itself, several real GPUs (one GPU here; NCCL refuses two processes on one
GPU, so the host path ran), or the speed of the exchanges (on one shared GPU each exchange incurs a
context switch: HMC with the multigrid 48\,s on two processes against 1.6\,s on one, which says
nothing about several GPUs).

\end{document}
